% TOP_PROBLEM: {"id": 2555, "title": "Kourovka Notebook Problem 21.46", "category_id": 17, "category": {"id": 17, "name": "group_theory", "display_name": "Group Theory", "description": "Problems about groups, group actions, representations, and related algebraic structures.", "slug": "group-theory", "order_index": 17, "created_at": "2026-07-31T15:26:25.671Z"}, "status": "open", "problem_number": "KOU-21.46", "set_id": 10}
% FIRST_POSTED: 2026-10-01
% ENTRY_KIND: statement_only
% UnsolvedMath ID 2555; source catalogue code KOU-21.46
% !TeX program = lualatex
% Definition and sources only; this file is not an LLM solution attempt.
\documentclass[11pt,a4paper]{article}
\usepackage[margin=25mm]{geometry}
\usepackage{fontspec}
\setmainfont{lmroman10-regular.otf}[BoldFont=lmroman10-bold.otf,
 ItalicFont=lmroman10-italic.otf,BoldItalicFont=lmroman10-bolditalic.otf]
\usepackage{amsmath,amssymb,mathtools}
\usepackage{unicode-math}
\setmathfont{latinmodern-math.otf}
\AtBeginDocument{\let\setminus\smallsetminus}
\usepackage{xurl,hyperref}
\hypersetup{colorlinks=true,urlcolor=blue}
\setcounter{secnumdepth}{0}
\setlength{\parindent}{0pt}
\setlength{\parskip}{5pt}
\setlength{\emergencystretch}{3em}
\begin{document}
\section*{Kourovka Notebook Problem 21.46}\label{2555}
{\small UnsolvedMath ID 2555 · KOU-21.46 · Group Theory · Kourovka Notebook - New Problems, Issue 21}
\subsection{Definitions and mathematical statement}
For a group $G$, its integral group ring $\mathbb ZG$ consists of
finite formal integer combinations of group elements. Give $\mathbb Z$
the trivial $\mathbb ZG$-module structure, so every $g\in G$ acts as
the identity. The ordinary integral cohomological dimension
$\operatorname{cd}_{\mathbb Z}(G)$ is the least length $d$ of a projective
resolution
\[
 0\longrightarrow P_d\longrightarrow\cdots\longrightarrow P_0
          \longrightarrow\mathbb Z\longrightarrow0
\]
over $\mathbb ZG$, or $\infty$ if no finite-length resolution exists.
Projective means a direct summand of a free module; finite generation
of the modules is not built into this definition.

The question asks whether there exists an infinite simple finitely
presented group $G$ with
\[
 3\le\operatorname{cd}_{\mathbb Z}(G)<\infty.
\]
Simple means nontrivial with no proper nontrivial normal subgroup.
Finitely presented means a quotient of a finitely generated free group
by the normal closure of finitely many relators. All three group-theoretic
requirements and the cohomological bound apply to the same group.

The dimension sought is an integer strictly greater than two, without
a specified upper value. A group of infinite cohomological dimension
is not a witness. Neither rational cohomological dimension nor virtual
cohomological dimension is substituted for the ordinary integral
invariant. A resolution gives an upper bound; to establish a qualifying
example one must also exclude dimensions zero, one and two, as well
as verify the finite presentation and simplicity.
\subsection{Short English statement}
Does there exist a finitely presented infinite simple group of finite integral cohomological dimension greater than two?
\subsection{Sources}
\begin{itemize}
\item[S1] ULAM.AI and the UnsolvedMath Contributors, supplied problems.json, record 2555 (KOU-21.46), Kourovka Notebook - New Problems, Issue 21. Catalogue identity and source transcription; CC BY 4.0.
\url{https://huggingface.co/datasets/ulamai/UnsolvedMath}.
\item[S2] The Kourovka Notebook, 21st edition, new problems, Problem 21.46. Expanded formulation of the supplied catalogue statement.
\url{https://arxiv.org/pdf/1401.0300v47}.
\end{itemize}
\end{document}
